\begin{examplebox}
\textbf{例 1【直接秒杀：焦点距离积】}

椭圆 $\frac{x^2}{25}+\frac{y^2}{9}=1$，求两焦点到\textbf{任意一条切线}的距离积。

\textbf{最短解：}由 $b^2=9$，立即得 $\boxed{d_1d_2=9}$。不需要设切点，也不需要求切线方程。

\medskip
\textbf{例 2【具体数值核验：斜率为 1】}

在例 1 的椭圆中，取切线 $l:y=x+\sqrt{34}$。因 $a^2k^2+b^2=25+9=34$，该直线确为切线。焦点为 $F_1(-4,0)$、$F_2(4,0)$，故
\[
d_1=\frac{\sqrt{34}-4}{\sqrt2},\qquad
d_2=\frac{\sqrt{34}+4}{\sqrt2},
\]
\[
d_1d_2=\frac{34-16}{2}=\boxed{9}.
\]
\textbf{观察：}两个距离不相等，但它们的积不会变。

\medskip
\textbf{例 3【原点距已知：无需知道切线方程】}

椭圆 $\frac{x^2}{25}+\frac{y^2}{9}=1$ 的一条切线到原点的距离为 $4$。求两焦点到该切线的距离。

\textbf{最快解：}已知 $d_1+d_2=2d_0=8$，$d_1d_2=b^2=9$，所以 $d_1,d_2$ 是
\[
t^2-8t+9=0
\]
的两根：
\[
\boxed{\{d_1,d_2\}=\{4-\sqrt7,\ 4+\sqrt7\}}.
\]
无需确定切线在椭圆的哪一侧，也无须求斜率。

\medskip
\textbf{例 4【已知垂足：以圆代椭圆】}

已知椭圆 $\frac{x^2}{25}+\frac{y^2}{16}=1$，右焦点为 $F_2(3,0)$。一条切线 $l$ 上，从 $F_2$ 引垂线的垂足为 $H(0,5)$。求该切线。

\textbf{分析：}先验算 $H$ 在主圆上：$0^2+5^2=25$。因 $F_2H=(-3,5)$ 与 $l$ 垂直，故 $l$ 的法向量取 $(-3,5)$，再代入 $H$：
\[
-3(x-0)+5(y-5)=0
\quad\Longrightarrow\quad
\boxed{-3x+5y=25}.
\]
\textbf{核验切线：}$25^2=25\cdot(-3)^2+16\cdot5^2=225+400=625$，确实相切。

\medskip
\textbf{例 5【设点法：给定非轴端点切线】}

椭圆 $\frac{x^2}{25}+\frac{y^2}{16}=1$ 上的点 $P(3,\frac{16}{5})$，求切线及两焦点到切线的距离，验证主圆性质。

切线由“点换元”得
\[
\frac{3x}{25}+\frac{(16/5)y}{16}=1
\quad\Longrightarrow\quad\boxed{3x+5y=25}.
\]
焦点 $F_1(-3,0)$、$F_2(3,0)$，故
\[
d_1=\frac{25+9}{\sqrt{34}}=\sqrt{34},\qquad
d_2=\frac{25-9}{\sqrt{34}}=\frac{16}{\sqrt{34}}.
\]
于是 $\boxed{d_1d_2=16=b^2}$。右焦点向该切线的垂足为
\[
H_2=(3,0)+\frac{16}{34}(3,5)
=\left(\frac{75}{17},\frac{40}{17}\right).
\]
代入主圆：
\[
OH_2^2=\frac{75^2+40^2}{17^2}=\frac{7225}{289}=25=a^2.
\]
\textbf{提示：}切点是 $P=(3,\frac{16}{5})$，垂足是 $H_2=(\frac{75}{17},\frac{40}{17})$，两者不同。

\medskip
\textbf{例 6【补充条件才可唯一确定椭圆】}

某椭圆中心在原点、焦点在 $x$ 轴上，两焦点到任意切线的距离积为 $4$，且椭圆经过 $(0,2)$。求椭圆。

由 $d_1d_2=b^2=4$，只知 $b=2$；而 $(0,2)$ 是所有形如
\[
\frac{x^2}{a^2}+\frac{y^2}{4}=1,\qquad a>2
\]
的椭圆的上顶点。\textbf{原题信息不足，不能唯一求出 $a$。}

若\textbf{再增加}条件“焦距为 $2\sqrt5$”（即 $c=\sqrt5$），才有 $a^2=b^2+c^2=9$，从而唯一得到
\[
\boxed{\frac{x^2}{9}+\frac{y^2}{4}=1}.
\]

\medskip
\textbf{例 7【极端位置检查：水平与竖直切线】}

对 $\frac{x^2}{25}+\frac{y^2}{9}=1$：当 $l:y=3$ 时，$d_1=d_2=3$，积为 $9$，中心距 $3=b$；当 $l:x=5$ 时，$\{d_1,d_2\}=\{9,1\}$，积仍为 $9$，中心距 $5=a$。这个例子同时检查\textbf{斜率为零}和\textbf{斜率不存在}的两种边界。
\end{examplebox}
